\documentclass{article}
\usepackage{amsmath}
\usepackage{array}
\usepackage{booktabs}
\usepackage{amssymb}
\usepackage{braket}

\begin{document}

\section*{Collocation Method: Scaling to Three Qubits}

\subsection*{1. Setup}

We return to the boundary value problem from Experiment~2:
\begin{equation}
    u''(t) + u'(t) + u(t) = t^2 + 3t + 3, \qquad t \in [-1, 1],
\end{equation}
with boundary conditions $u(-1) = 0$, $u(1) = 2$ and exact solution $u(t) = t^2 + t$.

In Experiments~1 and~2 we used $N = 4$ (two qubits): two boundary rows and two
interior collocation rows, with a degree-3 polynomial trial function.  Here we
scale to $N = 8$ (three qubits): two boundary rows and \textbf{six} interior
collocation rows, with a degree-7 polynomial trial function
\begin{equation}
    u(t) = \sum_{j=0}^{7} a_j\,\phi_j(t).
\end{equation}
The six equidistant interior collocation points on $(-1,1)$ are
\begin{equation}
    t_k = -1 + \frac{k}{7}, \qquad k = 1, 2, \ldots, 6,
\end{equation}
i.e.\ $t_1 \approx -0.714$, $t_2 \approx -0.429$, $t_3 \approx -0.143$,
$t_4 \approx 0.143$, $t_5 \approx 0.429$, $t_6 \approx 0.714$.

\subsection*{2. Case A: Monomial Basis}

Using the standard basis $\phi_j(t) = t^j$, the collocation matrix
$A_{\mathrm{mono}}$ is $8\times 8$ with
\begin{itemize}
    \item Row 0: evaluate $\phi_j(-1)$ (BC at $-1$),
    \item Row 1: evaluate $\phi_j(1)$ (BC at $+1$),
    \item Rows 2--7: evaluate $L[\phi_j](t_k)$ at each interior point,
\end{itemize}
where $L[\,\cdot\,] = (\,\cdot\,)'' + (\,\cdot\,)' + (\,\cdot\,)$.  The matrix
is $8\times 8$, non-singular ($\det A_{\mathrm{mono}} \approx 1003.4$), and
decomposes into all $4^3 = \mathbf{64}$ three-qubit Pauli strings.  Solving the
system recovers coefficients $a_0=0$, $a_1=1$, $a_2=1$, $a_j=0$ for $j\geq 3$,
confirming $u(t)=t^2+t$.

\subsection*{3. Pauli Count Explosion}

Comparing the three experiments:

\begin{center}
\begin{tabular}{cccccc}
\toprule
Experiment & Operator & $N$ & Qubits & Pauli strings (Case A) & Total possible \\
\midrule
1 & $u''=0$                   & 4 & 2 & 16 & 16 \\
2 & $u''+u'+u=f$              & 4 & 2 & 16 & 16 \\
3 & $u''+u'+u=f$ (scaled)     & 8 & 3 & 64 & 64 \\
\bottomrule
\end{tabular}
\end{center}

The monomial basis saturates all available Pauli strings in every case: adding
more collocation points does not improve the Pauli structure---it only makes it
exponentially worse.

\subsection*{4. Global Optimum via Linear Programming}

\paragraph{LP formulation.}
By the basis-change identity $M = A_{\mathrm{mono}}\,B^{\top}$ (established in
Experiments~1 and~2), every non-singular collocation matrix $M$ is reachable by
choosing $B = (A_{\mathrm{mono}}^{-1}\,M)^{\top}$.  Each Pauli coefficient
$c_P = \tfrac{1}{8}\operatorname{Tr}(P\,M)$ is linear in the entries of $M$, so
we solve
\begin{equation}
    \min_{M \in \mathbb{R}^{8\times 8}}\;\sum_{P}|c_P(M)|
    \quad\text{s.t.}\quad \operatorname{Tr}(M) = 8.
\end{equation}
The problem has $64 + 64 = 128$ variables (matrix entries plus slacks) and
$128 + 1$ constraints.

\paragraph{Result.}
The LP (HiGHS solver) terminates with status \emph{Optimal} and finds
\begin{equation}
    M^* = I_8,
    \qquad
    c_{III} = 1, \quad c_P = 0 \;\text{ for all other } P.
\end{equation}
This is a single Pauli string out of 64 — a \textbf{98.4\% reduction} from
Case~A.

\paragraph{Corresponding basis.}
The dual basis $B = A_{\mathrm{mono}}^{-\top}$ (numerically, $\det B \approx
9.97\times 10^{-4} \neq 0$) has coefficients:
\begin{align*}
\phi_0(t) &\approx 0.4681 - 0.5788\,t + 0.0554\,t^2 + 0.0780\,t^3
            - 0.0241\,t^4 + 0.0009\,t^5 + 0.0007\,t^6 - 0.0001\,t^7,\\
\phi_1(t) &\approx 1.2724 + 0.3010\,t - 0.7867\,t^2 + 0.2121\,t^3
            + 0.0126\,t^4 - 0.0131\,t^5 + 0.0017\,t^6 + 0.0001\,t^7,\\
\phi_2(t) &\approx -0.0521 + 0.0775\,t - 0.0062\,t^2 - 0.0136\,t^3
            - 0.0558\,t^4 + 0.0558\,t^5 + 0.1140\,t^6 - 0.1198\,t^7,\\
\phi_3(t) &\approx -0.1094 + 0.0599\,t - 0.0276\,t^2 + 0.0370\,t^3
            + 0.4409\,t^4 - 0.6676\,t^5 - 0.3039\,t^6 + 0.5707\,t^7,
\end{align*}
and $\phi_4,\ldots,\phi_7$ similarly (higher-degree corrections to the remaining
collocation rows).  Solving $I_8\,\mathbf{a}=\mathbf{b}$ gives $\mathbf{a}=\mathbf{b}$
directly.  The solution is verified numerically at five test points:

\begin{center}
\begin{tabular}{ccc}
\toprule
$t$ & $u_{\mathrm{computed}}$ & $u_{\mathrm{exact}} = t^2+t$ \\
\midrule
$-1.0$ & $0.000\,000\,00$ & $0.000\,000\,00$ \\
$-0.5$ & $-0.250\,000\,00$ & $-0.250\,000\,00$ \\
$\phantom{-}0.0$ & $0.000\,000\,00$ & $0.000\,000\,00$ \\
$\phantom{-}0.5$ & $0.750\,000\,00$ & $0.750\,000\,00$ \\
$\phantom{-}1.0$ & $2.000\,000\,00$ & $2.000\,000\,00$ \\
\bottomrule
\end{tabular}
\end{center}

\subsection*{5. The Universal Pattern}

Across all three experiments the LP always returns $M^* = I_N$:

\begin{center}
\begin{tabular}{ccccc}
\toprule
Experiment & $N$ & Qubits & Case A & LP optimum \\
\midrule
1 & 4 & 2 & 16/16 & 1/16 \\
2 & 4 & 2 & 16/16 & 1/16 \\
3 & 8 & 3 & 64/64 & 1/64 \\
\bottomrule
\end{tabular}
\end{center}

\paragraph{Why $M^* = I_N$ is always optimal.}
The identity matrix minimises $\sum_P |c_P|$ subject to $\operatorname{Tr}(M)=N$
because $I_N$ is the unique matrix (up to orthogonal equivalence) that concentrates
all its Pauli weight into the identity string $I^{\otimes n}$.  Formally:
$c_{I\cdots I}(I_N) = \operatorname{Tr}(I_N^{\otimes n}\cdot I_N)/N = 1$ and
$c_P(I_N) = \operatorname{Tr}(P\cdot I_N)/N = \operatorname{Tr}(P)/N = 0$ for all
non-identity $P$, since every non-identity Pauli string is traceless.

\paragraph{Consequence.}
The globally optimal Pauli count is \textbf{always 1}, regardless of the
differential operator, domain, boundary conditions, or number of collocation
points.  The operator-specific information is entirely absorbed by the dual
basis $B$, which grows in complexity as $N$ increases but leaves the quantum
representation ($M^* = I_N$, one Pauli string) unchanged.

\paragraph{Implication for VQLS.}
With $M^* = I_N$ the VQLS cost function reduces to
$\|A\ket{\psi} - \ket{b}\|^2 / \|\ket{b}\|^2$ with $A = I$, so the variational
circuit has nothing to optimise — $\ket{\psi} = \ket{b}$ is trivially the solution.
This means $M^* = I_N$ is classically trivial and does not provide a quantum
advantage.  The open question is whether there exists a \emph{non-trivial}
class of ODEs or boundary conditions for which the LP optimum is greater than 1
but still polynomially small in the number of qubits, which would be the regime
where VQLS could offer a genuine speedup.

\end{document}
